\scanpage{178}
Les \deleted{\unsure{automorphismes}} éléments de $\Omega(k)$ correspondent
ainsi aux systèmes $(A,(\unsure{B_i^*|B_i})_{i\geq1})$ d'éléments
de $\widehat S_k\{\{X_0\}\}$ \deleted{\unsure{de}},
\[
 \begin{array}{l|l}
 A=\displaystyle\sum_{j\geq1}a_jX_0^{(j)}&a_j,b_{ij}\in\widehat S_k=k\{\{T\}\}\\
 B_i=\displaystyle\sum_{j\geq0}b_{ij}X_0^{(j)}&
 \end{array}
\]
tels que
\[
 \begin{aligned}
 \text{a)}\quad&\epsilon(a_1)\;(=\Delta(\varphi_*))\in k^*\\
 &(\epsilon:\widehat S_k\mathbin{\text{---}}k\text{ l'augmentation})\\
 \text{b)}\quad&2B_i=0\quad\inserted{\forall i\geq1},\quad\text{i.e. }2b_{ij}=0
 \quad\forall i\geq1,j\geq0.
 \end{aligned}
\]
On peut aussi l'écrire comme l'ens. des systèmes
$(a_{jh})_{j\geq1,h\geq0}$ $(b_{ijh})_{i\geq1,j\geq0,h\geq0}$
d'éléments de $k$ (qui \deleted{sont les coeff.} dans les
\inserted{\unsure{développements de}}\deleted{\unsure{aux}}
$a_j$ et $b_{ij}$ comme séries à p.d. en $T$) avec les seules conditions
\[
 a_{10}\in k^*,\quad 2b_{ijh}=0.
\]
Pour $k$ variable, \deleted{l'hom}\inserted{le foncteur} $\Omega(k)$ est donc
représentable par \inserted{le schéma affine sur $\mathbb Z$}
\[
 G_m\mathbin{\unsure{\times_{\mathbb Z}|\times}}
 {}_2G_{\unsure{o|a}}^{\mathbb N^*\times\mathbb N}\times_{\mathbb Z}{}_2G_a^{\mathbb N^*\times\mathbb N\times\mathbb N}
\]
\[
 \Omega\simeq\operatorname{Spec}\left(
 \mathbb Z[\deleted{\unsure{T}}(A_{jh})_{j\geq1,h\geq0},
 (B_{ijh})_{i\geq1,j\geq0,h\geq0}]
 \right.
\]
modulo idéal engendré par les $2B_{ijh}$ et localisé de plus p.r. à $A_{1,0}$

\inserted{Ainsi} $\Omega$ peut être considéré comme un schéma
en groupes sur $\mathbb Z$, qui est le schéma.
des autom. du foncteur des $\mathbb Z$-algèbres
vers les $k$-alg. différentielles $\unsure{\tfrac12}$-simpliciales
$k\longmapsto\widehat{\mathcal A}_{k*}$. (On \unsure{envisage} comme \unsure{on|mor}
\margin{177}
